% ======================================================================
% THESIS DRAFT -- edit this file for metadata, prose and chapter order.
% Compile from thesis_report/ with XeLaTeX (latexmk -xelatex main.tex).
% Overleaf: upload this folder, select main.tex, choose XeLaTeX.
% Formatting: preamble.tex | References: references.bib | Help: README.md
% All chapter text stays here so sections are easy to exchange with an AI.
% ======================================================================
\documentclass[11pt,a4paper,oneside,openany]{memoir}
% Formatting lives here; edit thesis prose and metadata in main.tex.
% Based on the supplied supplements/example.tex and its unchanged KU package.
\usepackage{fontspec}
\usepackage{libertine}
\usepackage[english]{babel}
\usepackage{amsmath}
\usepackage{booktabs}
\usepackage{tabularx}
\usepackage{microtype}
\usepackage{xcolor}
\usepackage{graphicx}
\usepackage{pdflscape}
\usepackage{csquotes}
\usepackage[backend=biber,style=authoryear,sorting=nyt,maxcitenames=2]{biblatex}
\addbibresource{references.bib}

% Find the original KU package and its logos without moving or changing them.
\makeatletter
\def\input@path{{supplements/}}
\makeatother
\graphicspath{{figures/}{supplements/}}
\usepackage[english,samf,hyperref]{ku-frontpage}
% The template requests a semibold series absent from this font setup.
% Select bold explicitly instead of relying on repeated font substitutions.
\renewcommand{\KUsemibold}{\bfseries}
\pdfstringdefDisableCommands{\def\\{ }}

% Memoir controls the body layout; the KU package controls the cover.
\setlrmarginsandblock{30mm}{30mm}{*}
\setulmarginsandblock{28mm}{28mm}{*}
\checkandfixthelayout
\OnehalfSpacing
\chapterstyle{hangnum}
\pagestyle{plain}
\setcounter{secnumdepth}{2}
\setcounter{tocdepth}{1}
\setlength{\emergencystretch}{2em}

% Memoir can prevent the template from patching some title commands.
% These fallbacks support both patched and unpatched versions of memoir.
\providecommand{\frontpageauthor}[1]{}
\providecommand{\frontpagetitle}[1]{}
\providecommand{\frontpagesubtitle}[1]{}
\providecommand{\frontpagedate}[1]{}

% Visible writing prompts can be hidden using \draftnotesfalse in main.tex.
\newif\ifdraftnotes
\draftnotestrue
\newcommand{\draftnote}[1]{%
  \ifdraftnotes
    \par\smallskip
    {\small\noindent\textcolor{KU}{\textbf{Draft task:} #1}\par}%
    \smallskip
  \fi
}


% ---- EDIT: cover details -----------------------------------------------
% Working title: revise freely. Keep it short to fit the example's cover.
\newcommand{\ThesisTitle}{Urban Accessibility\\and Airbnb Prices}
\newcommand{\ThesisSubtitle}{Euclidean and walking measures across three Nordic capitals}
\newcommand{\ThesisAuthor}{Ahmad Jibril Hemdi}
\newcommand{\ThesisDate}{Submission date to be confirmed}
\newcommand{\ThesisAdvisors}{Supervisors: Sandro Sousa \& Jason Burton}

\assignment{Master's thesis}
\author{\ThesisAuthor}
\title{\ThesisTitle}
\subtitle{\ThesisSubtitle}
\date{\ThesisDate}
\advisor{\ThesisAdvisors}
\frontpageauthor{\ThesisAuthor}
\frontpagetitle{\ThesisTitle}
\frontpagesubtitle{\ThesisSubtitle}
\frontpagedate{\ThesisDate}
% Optional cover illustration, placed in figures/:
% \frontpageimage{cover.png}
\kupdfsetup{\ThesisTitle}{MSc Social Data Science}{\ThesisAuthor}

% ---- EDIT: writing mode -----------------------------------------------
\draftnotestrue % Change to \draftnotesfalse to hide visible writing prompts.
% Hiding prompts does not mean the draft is complete; search for them first.

\begin{document}
\frontmatter
\begingroup
  \SingleSpacing
  \maketitle
\endgroup

% ======================================================================
% FRONT MATTER -- write the abstract after the results and conclusion.
% ======================================================================
\chapter*{Abstract}
\addcontentsline{toc}{chapter}{Abstract}
\draftnote{Write the final abstract last: problem, three-city data, matched
measurements, geographic evaluation, main qualified findings and contribution.
Report asking prices and observational evidence. Add verified numbers only.}

% Add acknowledgements or required declarations here after checking the
% programme's requirements. The example's PhD statement is not used.

\cleardoublepage
\tableofcontents
% Uncomment once actual figures/tables have been inserted:
% \listoffigures
% \listoftables

\mainmatter
% ======================================================================
% CHAPTER 1 -- INTRODUCTION
% Purpose: establish the problem, the gap, the questions and the contribution.
% Authority: ../docs/design/scientific_proposal.md (not a literature review).
% ======================================================================
\chapter{Introduction}
\label{ch:introduction}

\section{Motivation and research problem}
\label{sec:motivation}
This thesis examines how straight-line and walking-network representations of
urban accessibility relate to Airbnb nightly asking prices in
Copenhagen--Frederiksberg, Oslo and Stockholm. It compares what the two
representations measure, the information they add to price prediction,
geographic variation within cities, and transfer between cities.

\draftnote{Explain why accessibility measurement matters for the research
problem. Define asking price early. Support the motivation with literature
you have read and verified; distinguish listings from bookings or revenue.}

\section{Research gap and contribution}
\label{sec:gap}
\draftnote{Build one argument from the literature to the gap: representation,
geographic evaluation and transfer. State the contribution of their combined
comparison without claiming universal novelty or a causal effect.}

\section{Research questions}
\label{sec:research-questions}
The analysis is organized around four research questions:
\begin{enumerate}
  \item[RQ1.] How do matched Euclidean and walking accessibility measurements
  differ across the three cities?
  \item[RQ2.] What information do these measurements add to within-city price
  prediction, and what complementary adjusted associations are observed?
  \item[RQ3.] How do predictive performance and accessibility contributions
  vary across geographic validation designs and held-out neighbourhoods?
  \item[RQ4.] How do relative-price predictions and accessibility contributions
  transfer across the six source--target city directions?
\end{enumerate}

\section{Scope and thesis roadmap}
\label{sec:roadmap}
The study uses June 2026 Inside Airbnb listings for conventional entire
dwellings. Its evidence is observational. Chapter~\ref{ch:literature} develops
the conceptual framework; Chapter~\ref{ch:data} describes the data and study
areas; and Chapter~\ref{ch:methods} defines measurement and evaluation.
Chapter~\ref{ch:results} presents the evidence in research-question order,
followed by discussion in Chapter~\ref{ch:discussion} and conclusions in
Chapter~\ref{ch:conclusion}.

% ======================================================================
% CHAPTER 2 -- CONCEPTS AND LITERATURE
% Purpose: explain the reasoning behind the design before technical details.
% Add only checked sources to references.bib; use \textcite / \parencite.
% ======================================================================
\chapter{Conceptual Framework and Related Literature}
\label{ch:literature}

\section{Accessibility and the spatial context of asking prices}
\label{sec:accessibility-concept}
\draftnote{Define accessibility, proximity and reachable opportunities.
Explain plausible relationships to asking prices, then identify confounding
and selection issues that prevent a causal interpretation here.}

\section{Euclidean and walking representations}
\label{sec:representations-literature}
\draftnote{Explain straight-line distance, network restrictions and
threshold-based opportunities. Build an expectation for measurement
differences without assuming that a richer network measure predicts better.}

\section{Prediction, geographic generalization and transfer}
\label{sec:generalization-literature}
\draftnote{Connect price-modelling literature to spatial dependence,
geographic holdouts and source--target distribution differences. Distinguish
prediction, adjusted association and transfer as separate questions.}

\section{Synthesis and analytical expectations}
\label{sec:expectations}
\draftnote{End with a short bridge to RQ1--RQ4. Label expectations as
conceptual; do not retrospectively call observed patterns preregistered
hypotheses. Avoid another full summary of the introduction.}

% ======================================================================
% CHAPTER 3 -- STUDY AREAS AND DATA
% Purpose: explain what is observed and which populations the results cover.
% Sources: ../docs/analysis/data_dictionary.md, source_registry.md,
% sample_audit.json, context_audit.json; ../outputs/tables/sample_flow.md.
% ======================================================================
\chapter{Study Areas and Data}
\label{ch:data}

\section{Study areas and operational neighbourhoods}
\label{sec:study-areas}
\draftnote{Introduce Copenhagen--Frederiksberg, Oslo and Stockholm. Separate
accepted study masks, provider validation groups and native density areas.
Explain the 11/17/11 operational groups without certifying them as current
official boundaries. Insert study\_areas.png and retain its CRS/source notes.}

\section{Sources, dates and geographic coverage}
\label{sec:sources}
\draftnote{Summarize Inside Airbnb, historical OSM networks/destinations and
native official population density. Separate acquisition dates from reference
dates. State coverage, licensing and vintage limitations. Income and
residential prices are excluded. Put detailed provenance in the appendix.}

\section{Sample construction and missingness}
\label{sec:samples}
\draftnote{Walk from raw listings to eligible positive-price references,
common predictive samples and genuine complete-case association samples.
Use the generated sample-flow and city-summary tables. Explain missing-price
exclusion and training-only predictive preprocessing separately.}

\section{Outcome and control variables}
\label{sec:variables}
Within-city models use the natural logarithm of original observed positive
nightly asking prices in native currency. Cross-city transfer uses the
city-relative outcome defined in Section~\ref{sec:transfer-method}.
Missing outcomes are never imputed.

\draftnote{Describe the common dwelling, host, amenity, density and centrality
controls with units and transformations. Record the native density
denominators and vintage caveats, especially Stockholm. Place the complete
variable dictionary in Appendix~\ref{app:definitions}.}

% ======================================================================
% CHAPTER 4 -- METHODS
% Purpose: explain how each question is answered; keep findings in Chapter 5.
% Authorities: ../config/study/design.json; ../docs/analysis/
% preanalysis_specification.md; association_missingness_amendment_20261013.md.
% ======================================================================
\chapter{Methods}
\label{ch:methods}

\section{Analytical design and model comparisons}
\label{sec:model-design}
Let $X$ denote the common controls, $E$ Euclidean accessibility and $W$
walking accessibility. The four model specifications are
\[
  M_0=X,\qquad M_E=X+E,\qquad M_W=X+W,\qquad M_{EW}=X+E+W.
\]
Within each city, all four predictive models use the same observations,
outcomes and evaluation folds. OLS and XGBoost are the specified estimators.

\draftnote{Explain the baseline, accessibility-versus-baseline comparisons,
walking replacement ($M_W-M_E$) and walking add-on ($M_{EW}-M_E$).
Introduce a compact design diagram only if it makes the sequence clearer.}

\section{Constructing matched accessibility measurements}
\label{sec:accessibility-method}
\draftnote{Define nearest rail/bus access and food/social and culture/tourism
opportunities on matched category universes. Explain the historical OSM
timestamp, directed graph and both snap connectors. The primary comparison
uses 5 km/h, 10 minutes and a rounded matched 833 m radius. Nearest destinations
may differ between representations; do not label their excess same-object
route circuity. Distinguish measured zeros from routing/source missingness.}

\section{Geographic validation and training procedure}
\label{sec:validation-method}
\draftnote{Explain leave-one-operational-neighbourhood-out evaluation and
the fixed 15-test/500-training support gate. Then describe random five-fold
and fixed-grid comparisons. Explain training-only fitted preprocessing and
the 30-candidate, four-inner-spatial-fold XGBoost search. Outer test folds
never select variables, thresholds, preprocessing or hyperparameters.}

\section{Metrics, paired contrasts and stability}
\label{sec:metrics-method}
\draftnote{Define RMSE and supporting metrics on their outcome scales.
Define candidate-minus-reference contrasts so negative RMSE change means
improvement. Explain 2,000 group resamples as conditional
geographic-composition stability, with overlapping training data and few
groups. These ranges do not establish independent-sample inference or
equivalence.}

\section{Complementary complete-case associations}
\label{sec:association-method}
\draftnote{Separate adjusted OLS associations from predictive evaluation.
Explain common complete cases on $X+E+W$, HC3, restricted Webb inference and
spatial diagnostics/sensitivities. Disclose the uniform 5\% to 10\% eligibility
amendment on 13 October 2026 after observing Stockholm's missingness.
Missing bedrooms and outcomes remain unfilled in association inference.}

\section{Source-only cross-city transfer}
\label{sec:transfer-method}
For listing $i$ in city $c$, the transfer outcome is
\[
  y^{\mathrm{rel}}_{ic}=\log(p_{ic})-\mu_c,
\]
where $p_{ic}$ is observed positive nightly asking price and $\mu_c$ is the
frozen mean log price of the city's eligible inside-mask positive-price
reference population before feature filtering.

\draftnote{Describe all six directions. Preprocessing, tuning and fitting
use the source city only. Target prices supply the frozen target benchmark
and evaluation, with no target refitting or recalibration. Explain why this
is a relative-price transfer test and why directions are dependent.}

\section{Registered robustness and reproducibility}
\label{sec:robustness-method}
\draftnote{Group the completed checks by validation, sample/control,
accessibility and association purpose. Explain what each challenges rather
than repeating every setting. Keep full configurations and NOT\_APPLICABLE
statuses in Appendix~\ref{app:robustness}. Record the metric CRSs:
EPSG:25832 (Copenhagen), EPSG:25833 (Oslo), EPSG:3006 (Stockholm).
Describe code/SQL provenance without exposing private records.}

% ======================================================================
% CHAPTER 5 -- RESULTS (same order as the four research questions)
% Purpose: answer each question with evidence before discussing mechanisms.
% Primary handover: ../reports/findings_by_research_question.md,
% research_results.md and thesis_figure_recommendations.md.
% Every number/figure must trace to the pinned final generated outputs.
% ======================================================================
\chapter{Results}
\label{ch:results}

\section{RQ1: Accessibility measurement differences}
\label{sec:rq1-results}
% Evidence: accessibility_descriptive_statistics; nearest_access_divergence.
% Figures: nearest_access_relationship.png; opportunity_comparison.png.
\draftnote{Start with a direct answer. Compare nearest-category access and
opportunities with units and city differences. Insert the generated summaries,
then close with what this establishes for the modelling comparison.}

\section{RQ2: Within-city prediction and adjusted associations}
\label{sec:rq2-results}
\subsection{Predictive performance and accessibility contributions}
\label{sec:rq2-prediction}
% Tables: primary_predictive_performance; paired_model_contrasts.
% Figures: primary_model_performance.png; paired_accessibility_contrasts.png.
\draftnote{Show all four models and both estimators before emphasizing a
contrast. Separate baseline gains, replacement and add-on. The frozen
XGBoost point gains in Copenhagen/Stockholm and adverse Oslo changes require
city-specific qualifications and stability ranges.}

\subsection{Complementary adjusted associations}
\label{sec:rq2-associations}
% Evidence: association_results; appendix figure association_accessibility.png.
\draftnote{Summarize the complete-case association evidence with units and
uncertainty. Keep full coefficients/diagnostics in the appendix. Do not
equate coefficients or descriptive feature contributions with causal effects
or held-out predictive gains.}

\section{RQ3: Geographic generalization and heterogeneity}
\label{sec:rq3-results}
% Evidence: validation_comparison; neighbourhood_performance.
% Figures: validation_strategy_comparison.png;
% neighbourhood_performance_variation.png.
\draftnote{Compare primary, random and fixed-block tasks, then show
neighbourhood variation. Explain reversals and unequal support. Summarize
registered robustness patterns without ranking scores across changed samples
or counting overlapping grids as independent replications.}

\section{RQ4: Directional cross-city transfer}
\label{sec:rq4-results}
% Tables: transfer_performance; transfer_contrasts.
% Figures: transfer_performance.png; transfer_network_comparison.png;
% directional_transfer_asymmetry.png. Use landscape when needed.
\draftnote{Cover every direction and both estimators. Explain relative-price
errors, baseline gains, walking replacement/add-on and directional asymmetry.
Do not average away reversals or imply six independent replications.}

\section{Summary of answers}
\label{sec:results-summary}
\draftnote{Give one concise, evidence-qualified answer to each question.
Use this as the bridge to discussion; do not repeat all tables or numerical
results. Check every statement against the current claims document.}

% ======================================================================
% CHAPTER 6 -- DISCUSSION
% Purpose: explain what the combined findings mean, with literature support.
% Claim boundaries: ../docs/final_claims_for_thesis.md.
% ======================================================================
\chapter{Discussion}
\label{ch:discussion}

\section{What accessibility representation changes}
\label{sec:discussion-measurement}
\draftnote{Connect RQ1 measurement differences to RQ2 predictive patterns and
the literature. Explain why representational detail and predictive gain can
diverge. Distinguish observed evidence from plausible explanations.}

\section{Geography and the limits of transportability}
\label{sec:discussion-geography}
\draftnote{Read RQ3 and RQ4 together. Discuss validation geography,
source--target differences and directional dependence. Explain the scope of
generalization supported by three observed city samples.}

\section{Limitations and methodological decisions}
\label{sec:limitations}
\draftnote{Connect each limitation to its consequence: asking-price and
missing-outcome selection, displaced listing coordinates, OSM completeness,
native density/vintage mismatches, operational groups, few clusters, correlated
features and residual spatial dependence. Disclose decision timing and the
post-observation 10\% amendment. Transfer also requires observed target prices
for centring. Avoid unsupported causal housing/displacement claims.}

\section{Implications and future research}
\label{sec:implications}
\draftnote{State what a researcher learns about choosing accessibility
representations and evaluation geography. Tie future work to specific
limitations. No-density prediction, same-object circuity and median-centred
transfer are possible future analyses, not completed study results.}

% ======================================================================
% CHAPTER 7 -- CONCLUSION
% Purpose: leave a direct answer to the original problem; introduce no evidence.
% ======================================================================
\chapter{Conclusion}
\label{ch:conclusion}
\draftnote{Answer RQ1--RQ4 briefly, then state the combined contribution and
its scope. End on what the comparison establishes about accessibility,
geographic evaluation and transfer. Do not add new results or a universal
walking-superiority claim.}

% ======================================================================
% APPENDICES -- detailed evidence that supports, but interrupts, the argument.
% ======================================================================
\appendix
\chapter{Data Provenance and Variable Definitions}
\label{app:definitions}
\draftnote{Include source/reference/acquisition dates, geographic extents,
CRSs, checksums where available, sample-flow details, category definitions,
controls and qualified native density denominators. Use aggregate evidence.}

\chapter{Extended Results and Diagnostics}
\label{app:extended-results}
\draftnote{Place full associations, descriptive held-out contributions,
neighbourhood detail, transfer diagnostics and spatial residual checks here.
Retain source captions and privacy suppression for published maps/bins.}

\chapter{Robustness and Specification Decisions}
\label{app:robustness}
\draftnote{Include the complete generated registry, fixed-grid/settings
details and explicit non-applicable entries. Record decision chronology,
especially support gates and the association amendment. Historical drafts
and unexecuted suggestions are not current results.}

\backmatter
\chapter*{References}
\addcontentsline{toc}{chapter}{References}
\draftnote{Add checked entries to references.bib and cite them using
\texttt{\string\parencite\{key\}} or \texttt{\string\textcite\{key\}}.
The initial bibliography is intentionally empty; no citations are invented.}
\printbibliography[heading=none]

\end{document}
